% page 208 du fac-simile (image p-207 du scan)
% bloc 157.5 x 242.0 mm ; marge gauche 8.0 mm
\pgc{-12.700mm}{0.000mm}{
\l{\cel{3}2OO.}
\l{}
\l{granda. I. (ecepte pri la staturo homala). Qua transiras la}
\l{\cel{3}dimensioni ordinara (partikulare pri alteso o longeso). -}
\l{\cel{3}II. Qua transiras la quanto ordinara pri quanto o qualeso.}
\l{\cel{3}- III. Qua transiras la nivelo, la grado ordinara pri}
\l{\cel{3}la rango, la situeso sociala. -IV. Qua transiras la ni-}
\l{\cel{3}velo, la grado ordinara, pri la qualeso di la spirito}
\l{\cel{3}(mento) o di la kordio. - EFIS.}
\l{}
\l{\sou{grandioza}. Qua havas karaktero impozanta di grandeso.-DEFIRS}
\l{}
\l{\sou{granito}. (\sou{geol}.)Roko fairala, ek feldspato, mikao, quarco,}
\l{\cel{3}asemblita en maso kompakta. - DEFIRS.}
\l{}
\l{\sou{granivora}. Qua su nutras per cerealo-grani. - DEFIS.}
\l{}
\l{\sou{grantar}. (\sou{trans.,}\cel{1}ad). Donar o koncesar ulo kom favoro}
\l{\cel{3}(do demandita o dezirita) . - E.}
\l{}
\l{\sou{granulo}. (\sou{farmacio)}. Pilulo extreme mikra. - DEFIRS.}
\l{}
\l{\sou{granulomo}. (\sou{patol.}) Mikra tumoro globatra, irgequala : sifi-}
\l{\cel{3}lisala, tuberklosala, inflamala, e c.}
\l{}
\l{\sou{grapo}. (\sou{bot}.) Asemblajo de flori, de frukti, dispozita eta-}
\l{\cel{3}jatre sur pedunklo komuna. - FI.}
\l{}
\l{\sou{grapino.}\cel{2}I. (\sou{navig}). l. Abordo-hoko. 2. Ankro di shalupo kun}
\l{\cel{3}plura pinti hokatra. - II. (\sou{tekn.}) l. Instrumento por sepa-}
\l{\cel{3}rar, en la vitbero-presilo, la grani de la raflo. 2. Ins-}
\l{\cel{3}trumento por deprenar la sordidaji de la vitro fuzesanta.}
\l{\cel{3}3. Ferajo hokatra di la kamen-skrapisto por skrapar la fu-}
\l{\cel{3}ligino. 4. Hoko quan onu fixigas an sua pedi por klimar}
\l{\cel{3}sur la arbori. - EFI.}
\l{}
\l{\sou{graso}. I. Substanco oleoza, uncion-apta, dispersita en la ti-}
\l{\cel{3}suo celulara di la homo-korpo e di la animali. - II. Ta}
\l{\cel{3}substanco, deprenita de la korpo di la animali, ed uzata}
\l{\cel{3}por koquo od en la industrii. - EFIS.}
\l{}
\l{\sou{gratar}. (\sou{trans.})\cel{2}Frotar, skrapante la surfaco. - FI.}
\l{}
\l{\sou{gratifikar}. (trans.)\cel{2}Donacar pekunio ad ulu kom atesto pri}
\l{\cel{3}sua satisfaceso, ultre to quo debesas a lu po lua laboro.}
\l{\cel{3}DFIS.}
\l{}
\l{\sou{gratino}. I. Parto di ula disho, qua adheras e risolas sur la}
\l{\cel{3}parieto di la vazo lor koqueso. - II. Maniero koquar ula}
\l{\cel{3}dishi kovrita de pan-krusto raspita,e pose risoligita.-DEF.}
\l{}
\l{\sou{gratitudar}. (\sou{trans., pri, pro}). Atestar ke onu etiko-debas ad}
\l{\cel{3}ulu pro ke onu recevis de lu bonfaco od helpo.- EFIS.}
\l{}
\l{\sou{gratuita}. I. Quan onu profitas sen pagar. - II.\cel{1}\sou{Gratuita}\cel{1}:}
\l{\cel{3}sen demandar pago kompense. - DEFIRS.}
}
